\section{Conclusion}
\label{sec:conclusion}

\sys\ addresses a concrete gap: while RSI-Forge \cite{rsi-forge} demonstrates
that papers can be automatically converted into executable environments, and
\harbor\ \cite{harbor} provides standardized infrastructure for agent evaluation
with existing life-science task sets, \emph{no public work bridges these
components systematically}. We propose a pipeline and a four-stage validation
protocol that progresses from environment quality verification through
weight-free self-improvement and harness optimization to post-training with
held-out benchmarks. The protocol aligns with RSI-Forge's existing evidence
chain while adding the causal evidence---post-training effects on unseen
tasks---required to claim that generated environments improve model
capabilities, not merely provide improvement signals. The initial rollout
focuses on computational biology and bioinformatics, the subfields where
automation is most feasible and impact is most immediate.

Because the empirical stages are expensive, we have separated the claim space
explicitly: the pipeline, its verification gates, and the protocol
implementation are released and measured (Appendix~\ref{app:claims}), while
every empirical number remains registered as \emph{planned} with the control arm
it will be measured against. The ledger is the contract with the reader: no
number becomes a result until a script in the repository produces it.

\paragraph{Author contributions.} Conceptualization, methodology, software,
writing---original draft preparation.

\paragraph{Competing interests.} The authors declare no competing interests.

\paragraph{Data availability.} Pipeline code and generated \harbor\
environments will be released upon publication.
